\section{Scope and Limitations}

This paper presents a conceptual framework, not a mathematical proof or empirically testable theory. The mathematical notation serves as a tool for precision in expressing philosophical concepts, following the tradition of using formal methods in philosophy of mind. However, the mathematical formalism should be understood as provisional notation for conceptual clarity rather than rigorous mathematical definitions.

Several significant limitations characterize this work. Most fundamentally, the framework currently lacks specific, falsifiable predictions that could be tested experimentally. Developing such predictions remains a significant challenge that must be addressed before the framework can move beyond philosophical speculation. Additionally, while mathematical language is employed throughout, formal proofs or demonstrated mathematical consistency have not been provided. The notation serves primarily to enhance conceptual precision rather than establish rigorous mathematical foundations.

Furthermore, although the framework aims to suggests a possible interpretation of consciousness, it has not been systematically connected to detailed phenomenological analysis or neuroscientific data. The rich literature of phenomenology and the growing body of neuroscientific evidence about consciousness remain largely unintegrated with these theoretical proposals. Finally, the suggested connections to physics are highly speculative. Deriving known physical laws from this framework would require substantial theoretical development that has not been attempted here.

Despite these limitations, this framework offers value as a coherent philosophical position that could inspire future research. It provides a novel perspective on the mind-body problem and suggests new ways of thinking about the relationship between experience and physical reality. This work aims to contribute to ongoing philosophical discussions and potentially inform future scientific approaches to consciousness.
